% !TEX root = task-1-13-3.tex
% Editable task bank. Append new BankTask blocks at the end of the document.
\documentclass[12pt,a4paper]{article}
\usepackage{physics-textbook}
\usepackage{graphicx,needspace}
\usetikzlibrary{decorations.pathmorphing}
\newif\ifShowSolutions
\ifdefined\TutorEdition
  \ShowSolutionstrue
\else
  \ShowSolutionsfalse % Change to \ShowSolutionstrue to display solutions.
\fi
\newcounter{banktask}
\newsavebox{\BankTaskBox}
\newcommand{\BankTaskContent}[3]{%
  \noindent\textbf{\thebanktask.}\quad #2\par
  \ifShowSolutions
    \Needspace{50mm}% Keep the solution heading with its opening text or diagram.
    #3\par
    \ifdefined\TutorEdition
      \ifcsname TutorSource#1\endcsname
        \nopagebreak\smallskip{\small\csname TutorSource#1\endcsname\par}%
      \fi
    \fi
  \fi
}
\newcommand{\BankTask}[3]{%
  \refstepcounter{banktask}\par\medskip
  \sbox{\BankTaskBox}{\begin{minipage}{\linewidth}%
    \BankTaskContent{#1}{#2}{#3}\end{minipage}}%
  \ifdim\dimexpr\ht\BankTaskBox+\dp\BankTaskBox\relax>\dimexpr\textheight-20mm\relax
    \Needspace{35mm}\BankTaskContent{#1}{#2}{#3}%
  \else
    \Needspace{\dimexpr\ht\BankTaskBox+\dp\BankTaskBox\relax}%
    \noindent\usebox{\BankTaskBox}\par
  \fi
  \medskip
}
\begin{document}
\BookSubsection{1.13.3}{Гидравлический пресс}


% Source sdamgia-9496: https://phys-ege.sdamgia.ru/problem?id=9496
\BankTask{sdamgia-9496}{
Гидравлический пресс изготовлен с использованием двух вертикальных цилиндрических сообщающихся сосудов, заполненных жидкостью и закрытых легкими поршнями. Радиус большего поршня этого пресса превосходит радиус меньшего поршня в $5$ раз. На малый поршень положили груз массой $20 \,\text{кг}$, удерживая больший поршень неподвижным. Определите модуль силы давления жидкости на больший поршень. Атмосферным давлением пренебречь.
}{
\smallskip\noindent\textbf{Решение}\par
Запишем закон гидравлического пресса:

$S_{2}/S_{1} = F_{2}/F_{1},$

где $S= \pi R ^{2}$ — площадь поршня, $F$ — сила, действующая на соответствующий поршень.
 
На малый поршень действует сила тяжести груза, а на большой — сила, по модулю равная силе давления жидкости (согласно третьему закону Ньютона).
Найдем, чему равна сила давления жидкости на большой поршень:

$F_{2}= \frac{ S_{2}F_{1}}{ S_{1} } = \frac{ R_{2} ^{2} }{ R_{1} ^{2} } \cdot mg=25 \cdot 200=5000\,\text{Н}.$

Ответ: $5000$.
}

% Source sdamgia-9528: https://phys-ege.sdamgia.ru/problem?id=9528
\BankTask{sdamgia-9528}{
Гидравлический пресс изготовлен с использованием двух вертикальных цилиндрических сообщающихся сосудов, заполненных жидкостью и закрытых легкими поршнями. Если на малый поршень этого пресса положить груз массой $40 \,\text{кг}$, удерживая больший поршень неподвижным, то сила давления жидкости на больший поршень увеличится на 
 $900 \,\text{Н}$. Во сколько раз радиус большего поршня пресса превосходит радиус меньшего поршня? Атмосферным давлением пренебречь.
}{
\smallskip\noindent\textbf{Решение}\par
Запишем закон гидравлического пресса:

$S_{2}/S_{1} = F_{2}/F_{1},$

где $S= \pi R ^{2}$ — площадь поршня, $F$ — сила, действующая на соответствующий поршень.
На малый поршень действует сила тяжести груза, а на большой — сила, по модулю равная силе давления жидкости (согласно третьему закону Ньютона).
Найдем чему равно отношение радиусов поршней:

$\frac{ R_{2} ^{2} }{ R_{1} ^{2} } = \frac{ F_{2}}{ F_{1} } = \frac{ F_{2}}{ mg } = \frac{ 9}{ 4 } \Longleftrightarrow \frac{ R_{2}}{ R_{1} } = \frac{ 3}{ 2 } .$

Таким образом, радиус большого поршня в $1{,}5$ раза больше радиуса первого поршня.
 
Ответ: $1{,}5$.
}

\end{document}
